\subsection{Learning Paradigm and Research Objective}
\label{subsec:matching-dm-methods}

The objective is to estimate gallstone presence from demographic, clinical,
laboratory, and body-composition measurements. All 319 records have a known
binary outcome, making \textbf{supervised learning} the appropriate primary
paradigm: predictions can be compared with observed labels on held-out records.
\textbf{Unsupervised learning} can reveal structure without those labels, but
similarity between patients does not necessarily correspond to disease
status. It therefore supports exploration rather than providing the principal
answer to the research question.

\subsection{Comparison of Data-Mining Methods}
\label{subsec:dm-family-comparison}

The learning paradigm must be narrowed to a method whose output matches the
objective. Table~\ref{tab:dm-family-comparison} distinguishes this decision
from the choice of a particular algorithm.

\begin{table}[H]
\centering
\small
\caption{Suitability of data-mining methods for the gallstone-prediction objective.}
\label{tab:dm-family-comparison}
\begin{tabularx}{\textwidth}{@{}p{2.8cm}XX@{}}
\toprule
Method & Intended output & Suitability for this project \\
\midrule
\textbf{Classification} &
A known class and, where supported, its probability. &
Directly matches labelled gallstone presence/absence and permits evaluation
of discrimination, missed cases, and false alarms. Selected as the primary
method. \\
Continuous-value prediction (regression) &
A numerical quantity, such as a laboratory concentration or severity score. &
The target is a category, not a continuous measurement. This would answer a
different question. Classification is also predictive; the distinction here
concerns the type of output. \\
Segmentation (clustering) &
Groups of observations with similar predictor profiles. &
Useful for exploratory subgroups, but groups need not correspond to
gallstone status. Cluster membership alone is not a validated disease
prediction. \\
Association discovery &
Frequently co-occurring conditions or measurement patterns. &
Useful for hypothesis generation, but support and confidence do not establish
held-out discrimination or an acceptable false-negative rate. Outcome-directed
rules would still require supervised classification evaluation. \\
Anomaly detection &
Observations that differ from a reference pattern. &
Supports review of unusual measurements, but gallstone cases comprise roughly
half the sample and cannot be equated with anomalies. \\
\bottomrule
\end{tabularx}
\end{table}

Dimensionality reduction and feature selection are supporting preparation
methods, not substitutes for predicting the labelled outcome. Likewise,
forecasting future onset or modelling time to diagnosis would require
longitudinal or event-time information not used by the current workflow.
The available evidence supports classification of recorded gallstone status,
not a demonstrated forecast of future disease.

\subsection{Selected Method and Success Criteria}
\label{subsec:dm-method-decision}

The selected method is \textbf{supervised binary classification with
probability estimation}. It provides the required patient-level output and
directly supports the project benchmarks of ROC--AUC at least 0.80 and
gallstone-class sensitivity at least 0.80. Specificity, precision, F1,
stability, and interpretability are needed alongside these measures: predicting
every patient as positive could satisfy sensitivity while producing
unacceptable false alarms.

The source coding is 0 = gallstones present and 1 = absent. The Python
modelling target is recoded to 1 = present and 0 = absent, so
$\mathrm{FNR}=\mathrm{FN}/(\mathrm{FN}+\mathrm{TP})=1-\mathrm{sensitivity}$
measures missed gallstone cases. These are analytical project benchmarks,
not clinical deployment thresholds. The following section compares algorithms
within this chosen method and explains how the implemented models are built
and selected.
